% s-199-p199-m002; original PDF 199, 200
\Aphorism{s-199-p199-m002}{284}
Der leidenschaftliche Mensch, wie er uns im Leben und auf der Bühne gewöhnlich vorgeführt wird, ist ein Amphibium: er entlädt seine innere Fülle und Spannung \Sperr{zugleich} seelisch und körperlich, — damit mildert er die volle Wucht der Reaction, die, in Ein Strombette gedrängt, schlimmstenfalls einmal das ganze seelische Gebiet überfluthen und verwüsten könnte. So erleichtern wir uns einen heroischen Entschluss durch heftige und laute Gebärden; in Verzweiflungsanfällen wüthen wir gegen den eigenen Körper, und wenn ein „unbegreiflich holdes Sehnen“ in uns aufblüht, treibt es uns immer noch, wie Faust, durch Wald und Wiesen hinzu\Sperr{gehn}. Nur eine Kunst, wie die ältere griechische Tragödie, in der ein ungeheurer Zwang von Rhythmus und Ceremonie aufgeboten wird, um über das Naturalistische und Pathologische Herr zu werden, nur ein \Greek{φρενοτέκτων} \Greek{ἀνήρ} wie Aeschylus, der noch die flugsüchtigste Leidenschaft zum Gang auf fester Erde zu zwingen weiss (mit rossgleich schreitenden Worten, wie Aristophanes sagt), durfte es wagen, einen Menschen, und noch dazu einen Prometheus, in vollster \Sperr{körperlicher Unbeweglichkeit} vorzuführen. Ein modernes Kunstwerk, das den gleichen Versuch unternähme, würde zwei Schritt vor der Irrenklinik enden; psychische Energie, die nach aussen gehemmt ins Seelische zurückgestaut wird, führt zur Degeneration. Darum werden Verbrecher im Kerker nicht „besser“ —.
\par

